\section{Attention as a restricted receiver-interface realization}\label{sec:attention}

The state-indexed operator family of \Cref{sec:provenance-routing} does not imply attention. Attention is obtained by restricting the source-resolved component of $Q_r$ to a particular scalar, positive, row-anchored, source-local, additive, and finitely separable form. The effective-potential representation appears only after that routing row has been derived.

The assumptions below select this narrower regime rather than axiomatizing receiver-organized computation in general. Their distinct architectural work is exhibited by explicit neighboring constructions in the technical supplement; removing one commitment changes a typed component of the realization rather than making the surrounding factorization meaningless.

Let $R$ be a finite receiver set, let $S$ be a finite source set, and let $H\in X$ be the current represented state. Fix typed receiver and source feature maps
\[
x_r:X\to\mathcal X_R,
\qquad
y_i:X\to\mathcal X_S.
\]
For self-attention, $R=S$ and one usually has $x_i(H)=y_i(H)=h_i$. Keeping the two feature families distinct makes the normal form equally well typed for cross-attention and other receiver--source architectures.

\subsection{Commitments before the formula}

\begin{assumption}[Pairwise scalar evidence]\label{ass:pairwise-scalar}
For every admissible pair $(r,i)$, the routed interface uses a finite scalar evidence coordinate
\[
s_{ri}(H)\in\R.
\]
Hard exclusion is represented separately by a mask.
\end{assumption}

\begin{assumption}[Positive baseline and nonempty rows]\label{ass:positive-baseline}
There is a hard mask $M_{ri}\in\{0,1\}$ and a positive baseline factor $\kappa_{ri}>0$ for every pair. For each receiver,
\[
S_r:=\{i\in S:M_{ri}=1\}
\]
is nonempty. Values of $\kappa_{ri}$ off the admissible support are immaterial.
\end{assumption}

The single factor $\kappa_{ri}$ contains every state-independent positive pair term that is extensionally relevant to the row. It may be decomposed into schema-, instance-, or parameter-origin factors only when those factors are independently marked. A finite additive pair bias is equivalent to multiplying $\kappa_{ri}$ by its exponential and is not counted twice.

\begin{assumption}[Compositional positive link]\label{ass:link}
There is a continuous strictly increasing map
\[
\psi:\R\to(0,\infty)
\]
with
\[
\psi(a+b)=\psi(a)\psi(b).
\]
Thus independent additive evidence coordinates combine multiplicatively as positive raw support.
\end{assumption}

\begin{assumption}[Receiver-row anchoring]\label{ass:row-anchor}
For each receiver $r$, positive raw weights $(w_{ri})_{i\in S_r}$ are converted to coefficients $(\alpha_{ri})_{i\in S}$ such that:
\begin{enumerate}[label=(\roman*)]
\item $\alpha_{ri}=0$ for $i\notin S_r$;
\item $\alpha_{ri}>0$ for $i\in S_r$ and $\sum_{i\in S_r}\alpha_{ri}=1$;
\item within-row ratios are preserved:
\[
\frac{\alpha_{ri}}{\alpha_{rj}}
=
\frac{w_{ri}}{w_{rj}}
\qquad(i,j\in S_r).
\]
\end{enumerate}
\end{assumption}

\begin{assumption}[Source-local payload and additive aggregate]\label{ass:values-aggregation}
Each source has a payload map
\[
v_i:\mathcal X_S\to W,
\]
so its evaluated payload $v_i(y_i(H))$ factors through the marked source feature and is independent of the receiver. The transported exposure is
\[
E_r(H)=z_r(H):=\sum_{i\in S}\alpha_{ri}(H)v_i(y_i(H)).
\]
\end{assumption}

The commitment chain is
\[
\boxed{
Q_r
\rightsquigarrow
P_r
\rightsquigarrow
\bigl(s_{ri},\kappa_{ri},M_{ri}\bigr)_i
\rightsquigarrow
\alpha_r
\rightsquigarrow
z_r.
}
\]
More general receiver interfaces may use higher-order contributions, signed or operator-valued coefficients, pair-conditioned payloads, unnormalized flow, globally coupled plans, recurrent aggregation, or retained source-resolved memory.

\subsection{The positive link and row anchor}

\begin{lemma}[Exponential evidence link]\label{lem:exponential-link}
Under \Cref{ass:link}, there exists $\beta>0$ such that
\[
\psi(t)=e^{\beta t}.
\]
\end{lemma}

\begin{proof}
The multiplicative law gives $\psi(0)=1$. The continuous map $\varphi=\log\psi$ satisfies the additive Cauchy equation, so $\varphi(t)=\beta t$. Strict monotonicity implies $\beta>0$.
\end{proof}

Set $\tau:=\beta^{-1}>0$. This is a score scale controlling the evidence--dispersion tradeoff; no physical temperature claim follows.

\begin{lemma}[Uniqueness of receiver-row anchoring]\label{lem:row-anchor}
Let $(w_i)_{i\in A}$ be positive on a nonempty finite set. The unique positive unit-mass vector preserving all ratios $w_i/w_j$ is
\[
\alpha_i=\frac{w_i}{\sum_{j\in A}w_j}.
\]
\end{lemma}

\begin{proof}
Ratio preservation implies $\alpha_i=cw_i$ for one $c>0$. Unit mass determines $c=(\sum_jw_j)^{-1}$.
\end{proof}

Thus, on admissible pairs,
\[
w_{ri}(H)=\kappa_{ri}e^{s_{ri}(H)/\tau}
\]
and
\begin{equation}\label{eq:general-softmax-row}
\alpha_{ri}(H)
=
\frac{M_{ri}\kappa_{ri}e^{s_{ri}(H)/\tau}}
{\sum_{j\in S_r}\kappa_{rj}e^{s_{rj}(H)/\tau}}.
\end{equation}
Row normalization is therefore not a primitive truth of architecture. It follows from positivity, receiver-local unit mass, and preservation of relative evidence \citep{luce1959individual}.

\subsection{Shared finite separation rank}

Low rank of one realized score matrix is not enough to produce shared query and key maps. The commitment concerns the score family across admissible states.

\begin{definition}[Shared finite separation rank]\label{def:shared-separation-rank}
The contextual score family has shared separation rank at most $d$ when there exist functions
\[
f_1,\ldots,f_d:\mathcal X_R\to\R,
\qquad
g_1,\ldots,g_d:\mathcal X_S\to\R
\]
such that, for every admissible pair $(r,i)$ and represented state $H$,
\begin{equation}\label{eq:separable-score}
s_{ri}(H)
=
\sum_{\ell=1}^{d}f_\ell(x_r(H))g_\ell(y_i(H)).
\end{equation}
The same feature functions are used across receivers, sources, and states.
\end{definition}

\begin{proposition}[Query--key coordinates from finite separation rank]\label{prop:separable-bilinear}
A contextual score family has the form \eqref{eq:separable-score} if and only if there are maps
\[
q:\mathcal X_R\to\R^d,
\qquad
k:\mathcal X_S\to\R^d
\]
with
\[
s_{ri}(H)=\ip{q(x_r(H))}{k(y_i(H))}
\]
for every admissible pair and represented state.
\end{proposition}

\begin{proof}
Given \eqref{eq:separable-score}, set $q(x)=(f_\ell(x))_\ell$ and $k(y)=(g_\ell(y))_\ell$. Evaluating these maps at $x_r(H)$ and $y_i(H)$ gives the stated score. The converse follows by expanding the Euclidean inner product coordinatewise.
\end{proof}

The query--key coordinates are therefore derived coordinates on one shared separable score family. They are not semantic primitives and are not uniquely determined.

\subsection{Attention normal form}

\begin{theorem}[Single-head attention normal form]\label{thm:attention-normal-form}
Assume \Cref{ass:pairwise-scalar,ass:positive-baseline,ass:link,ass:row-anchor,ass:values-aggregation} and shared separation rank at most $d$. Then there are shared query and key maps $q,k$ such that
\begin{equation}\label{eq:attention-normal-form}
\boxed{
z_r(H)=\sum_{i\in S}\alpha_{ri}(H)v_i(y_i(H)),
}
\end{equation}
where
\begin{equation}\label{eq:attention-weights}
\boxed{
\alpha_{ri}(H)
=
\frac{
M_{ri}\kappa_{ri}
\exp\!\left(\ip{q(x_r(H))}{k(y_i(H))}/\tau\right)
}{
\sum_{j\in S_r}\kappa_{rj}
\exp\!\left(\ip{q(x_r(H))}{k(y_j(H))}/\tau\right)
}.
}
\end{equation}
\end{theorem}

\begin{proof}
The exponential-link lemma gives positive raw weights $\kappa_{ri}e^{s_{ri}(H)/\tau}$. The separation-rank proposition gives $s_{ri}(H)=\ip{q(x_r(H))}{k(y_i(H))}$. The row-anchor lemma gives \eqref{eq:attention-weights}, and source-local payloads with additive aggregation give \eqref{eq:attention-normal-form}.
\end{proof}

\begin{corollary}[Standard projected attention]\label{cor:standard-attention}
In self-attention, if $R=S$, $x_i(H)=y_i(H)=h_i$, and
\[
q(h)=W_Qh,
\qquad
k(h)=W_Kh,
\qquad
v_i(h)=W_Vh,
\]
and a finite supplied pair bias $b_{ri}$ is represented by $\kappa_{ri}=e^{b_{ri}/\tau}$ on admissible pairs, then \Cref{thm:attention-normal-form} is masked scaled dot-product attention after choosing the conventional score scale \citep{vaswani2017attention}.
\end{corollary}

A state-dependent pair bias should instead remain in $s_{ri}(H)$. The decomposition between $\kappa$ and $s$ is provenance-relative: it is meaningful when the architecture marks which term is state-independent and which is constructed from current state.

\begin{proposition}[Converse verification]\label{prop:attention-converse}
A standard masked single-head attention layer with at least one admissible source per receiver satisfies the assumptions above. Hence its receiver exposure has both the routed normal form and the effective-potential representation derived below.
\end{proposition}

\begin{proof}
The layer uses finite bilinear scores on admissible pairs, positive exponentiation, rowwise normalization, source-local values, and additive aggregation. A finite mask bias can be represented by $\kappa$; forbidden pairs are represented by $M=0$.
\end{proof}

The theorem and converse classify a declared regime. They do not imply that every state-conditioned interface, every graph constructor, or every operator field is attention.

\subsection{Effective potential after the routing law}\label{subsec:effective-potential}

Normalize the baseline on each admissible row:
\begin{equation}\label{eq:reference-measure}
\pi_{ri}:=\frac{\kappa_{ri}}{\sum_{j\in S_r}\kappa_{rj}}
\qquad(i\in S_r).
\end{equation}
Then $\pi_r$ is a reference probability measure on $S_r$.

\begin{definition}[Receiver-relative effective potential]\label{def:effective-potential}
For $i\in S_r$, define
\begin{equation}\label{eq:effective-potential}
\mathcal U_{ri}(H):=-s_{ri}(H)-\tau\log\pi_{ri}.
\end{equation}
For $i\notin S_r$, set $\mathcal U_{ri}(H)=+\infty$.
\end{definition}

Then
\begin{equation}\label{eq:gibbs-row}
\alpha_{ri}(H)
=
\frac{e^{-\mathcal U_{ri}(H)/\tau}}
{\sum_{j\in S_r}e^{-\mathcal U_{rj}(H)/\tau}}.
\end{equation}
The mask chooses the accessible source sector, the baseline supplies the state-independent reference landscape, and current-state evidence deforms that landscape. This is an exact coordinate on the positive scalar routing law, not the primitive definition of architecture.

\begin{proposition}[Marked multiplicative factors give additive potential terms]\label{prop:additive-potential}
Suppose the positive raw coupling has an independently marked factorization
\[
K_{ri}(H)=\prod_{a=1}^{m}K_{ri}^{(a)}(H),
\qquad K_{ri}^{(a)}(H)>0.
\]
Then, with $\mathcal U_{ri}^{(a)}=-\tau\log K_{ri}^{(a)}$,
\[
-\tau\log K_{ri}=\sum_{a=1}^{m}\mathcal U_{ri}^{(a)}.
\]
\end{proposition}

\begin{proof}
Apply the logarithm to the product.
\end{proof}

Thus schema-, instance-, parameter-, and state-origin factors may contribute additive potential terms when those factors are independently marked. The extensional row does not uniquely determine such a provenance decomposition.

\subsection{Receiver free energy}\label{subsec:free-energy}

For $p\in\Delta(S_r)$, define
\begin{equation}\label{eq:receiver-free-energy}
\mathcal F_r[p;H]
:=-\sum_{i\in S_r}p_i s_{ri}(H)
+\tau D_{\mathrm{KL}}(p\Vert\pi_r).
\end{equation}
Equivalently,
\[
\mathcal F_r[p;H]
=
\sum_ip_i\mathcal U_{ri}(H)-\tau H(p).
\]

\begin{theorem}[Receiver free-energy characterization]\label{thm:receiver-free-energy}
Let
\[
\mathcal Z_r(H):=\sum_{i\in S_r}\pi_{ri}e^{s_{ri}(H)/\tau}.
\]
For every $p\in\Delta(S_r)$,
\begin{equation}\label{eq:free-energy-gap}
\boxed{
\mathcal F_r[p;H]
=
-\tau\log\mathcal Z_r(H)
+\tau D_{\mathrm{KL}}\!\left(p\Vert\alpha_r(\cdot\mid H)\right).
}
\end{equation}
Consequently, $\alpha_r(\cdot\mid H)$ is the unique minimizer of $\mathcal F_r[\cdot;H]$.
\end{theorem}

\begin{proof}
For admissible $i$,
\[
\log\frac{\alpha_{ri}(H)}{\pi_{ri}}
=\frac{s_{ri}(H)}{\tau}-\log\mathcal Z_r(H).
\]
Substitution into $D_{\mathrm{KL}}(p\Vert\alpha_r)$ and rearrangement gives \eqref{eq:free-energy-gap}. Nonnegativity of KL divergence gives the unique minimizer.
\end{proof}

The theorem characterizes the already-derived row. It does not postulate physical equilibrium for the network. It states that, conditional on the current receiver, admissible source sector, baseline, score, and row-simplex constraint, softmax is the unique allocation balancing evidence against departure from the reference measure. Other ensembles---balanced transport, column capacities, unnormalized flow, signed allocation, or hard minimum selection---would define different architectures over related score data \citep{cuturi2013sinkhorn}.

\subsection{Exact source-carrier quotienting}\label{subsec:coarse-potential}

An exact quotient of the source carrier preserves a full pushed-forward allocation and is different from the moment aggregation studied in \Cref{sec:interface-refinement}.

Let $q:S_r\twoheadrightarrow C_r$ be a source quotient. Define
\begin{equation}\label{eq:coarse-kernel}
K_r^{\mathrm{eff}}(c\mid H)
:=\sum_{q(i)=c}e^{-\mathcal U_{ri}(H)/\tau}
\end{equation}
and
\begin{equation}\label{eq:coarse-potential}
\mathcal U_r^{\mathrm{eff}}(c;H)
:=-\tau\log K_r^{\mathrm{eff}}(c\mid H).
\end{equation}

\begin{theorem}[Effective potential under a source quotient]\label{thm:potential-of-mean-force}
If $\bar\alpha_r=q_\#\alpha_r$, then
\[
\boxed{
\bar\alpha_r(c\mid H)
=
\frac{e^{-\mathcal U_r^{\mathrm{eff}}(c;H)/\tau}}
{\sum_{c'\in C_r}e^{-\mathcal U_r^{\mathrm{eff}}(c';H)/\tau}}.
}
\]
The partition sum is preserved.
\end{theorem}

\begin{proof}
The numerator is the sum of the microscopic Gibbs numerators over the fiber $q^{-1}(c)$. Summing over all quotient classes partitions the original source sum.
\end{proof}

The induced potential is a log-sum-exp over eliminated sources. If a quotient class contains $n_c$ sources with equal potential $u$, then
\[
\mathcal U_r^{\mathrm{eff}}(c)=u-\tau\log n_c.
\]
Microscopic multiplicity and potential combine. By contrast, ordinary value aggregation
\[
(\alpha_r,(v_i)_i)\longmapsto\sum_i\alpha_{ri}v_i
\]
generally retains only one selected moment and does not preserve the pushed-forward allocation.

\subsection{Gauges, observables, and heads}

\begin{proposition}[Receiver potential gauge]\label{prop:potential-gauge}
Two finite potentials on the same admissible row induce the same normalized allocation if and only if they differ by a receiver-dependent additive constant.
\end{proposition}

\begin{proof}
A common additive shift produces a common multiplicative Gibbs factor that cancels. Conversely, equality of normalized rows makes $e^{-\mathcal U_i'/\tau}/e^{-\mathcal U_i/\tau}$ independent of $i$.
\end{proof}

\begin{lemma}[Query--key coordinate gauge]\label{lem:qk-gauge}
For $R\in GL(d)$, the transformations
\[
q'(x)=Rq(x),
\qquad
k'(y)=R^{-\mathsf T}k(y)
\]
preserve every score $\ip{q(x)}{k(y)}$.
\end{lemma}

The potential gauge changes a landscape coordinate without changing the row. The query--key gauge changes coordinates without changing the score. Neither representation is uniquely individuated by routing behavior.

The value map attaches an observable to each source:
\[
z_r(H)=\E_{i\sim\alpha_r(\cdot\mid H)}[v_i(y_i(H))].
\]
Values are not generally part of the potential. Two heads may use the same allocation with different observables, or different allocations whose selected moments coincide.

\begin{corollary}[Multihead routed decomposition]\label{cor:multihead}
A multihead output
\[
z_r
=
\sum_{m\in\mathcal H}O_m\left(\sum_i\alpha_{ri}^{(m)}v_i^{(m)}\right)
\]
is a parallel family of head-indexed receiver-interface refinements followed by linear head recombination.
\end{corollary}

The heads do not form one joint probability law unless an additional joint normalization is imposed.

\subsection{Commitment ablation}

\begin{table}[t]
\centering
\small
\begin{tabularx}{\textwidth}{@{}lX@{}}
\toprule
Commitment changed & Neighboring receiver-interface regime \\
\midrule
Pairwise scalar evidence & higher-order, hyperedge, tuple, vector-, matrix-, or operator-valued contributions \\
Positive support & signed, inhibitory, or cancellation-based routing \\
Receiver-row anchoring & columnwise, global, balanced, capacity-constrained, sparse, hard, or unnormalized allocation \\
Source-local payload & pair-conditioned or receiver-conditioned messages \\
Additive aggregate & recurrent aggregation, nonlinear set state, or retained source-resolved memory \\
Shared finite separation rank & direct pair network, nonseparable kernel, or state-specific score coordinates \\
Fixed admissible sector & state-generated hard barriers or persistent learned relation states \\
One-shot exposure & iterative retrieval, recurrent source access, or global planning \\
\bottomrule
\end{tabularx}
\caption{The attention theorem as a map from architectural commitments to neighboring interface forms.}
\label{tab:attention-ablation}
\end{table}

The theorem therefore exposes, rather than eliminates, architecture. The Transformer endogenizes receiver-relative allocation but continues to fix the source carrier, pairwise scalar interaction language, positivity, row-simplex ensemble, value transport form, aggregate statistic, and one-step schedule.
